% Correlated analytic representation of alpha: action, convergence, and
% intervals.
\subsubsection{Analytic representation of the alpha coordinate}
\label{subsec:alpha-lector-completo-rev08}

The preceding analysis of the jet requires the tail to be exhibited:
merely naming a complete reader is insufficient. Let \(\mathsf s\in\mathcal X_\alpha^{\mathrm{enr}}\). Before
decimal normalization, the same state already expresses the pre-coordinate
\begin{equation}
 a(\mathsf s):=p(\mathsf s)+e_0(\mathsf s)-f(\mathsf s)
                 -\kappa_{\mathrm{per}}(\mathsf s)\in I_\alpha,
 \label{eq:alpha-precoordenada-comun}
\end{equation}
with \(p=\pi_{\mathrm{HMT}}-3\),
\(e_0=e_{\mathrm{HMT}}-2\), and
\(f=\varphi_{\mathrm{HMT}}-1\). Equivalently,
\(a=\pi_{\mathrm{HMT}}+e_{\mathrm{HMT}}-
\varphi_{\mathrm{HMT}}-4-\kappa_{\mathrm{per}}\). Conventional
\(\alpha\) is not introduced here: \(p,e_0,f\) are the three regional
characters already generated, and \(\kappa_{\mathrm{per}}\) is the rational reading
of the terminal record. The integer closure applies the normalization in
\eqref{eq:alpha-recurrencia} to
\eqref{eq:alpha-precoordenada-comun}; the chart below expresses a compatible
analytic representation on the same pre-coordinate. Its dependency graph
receives \((p,e_0,f,\kappa_{\mathrm{per}})\) directly: it contains no
edge \(\alpha_A\to\lambda\), although the equality proved below allows
the two representation names to be identified subsequently.

Write \(q=1/729\), let \(F_{\mathsf s}\) be the ordered field generated by
the coordinates of \(\mathsf s\), and set
\(E_{9,\mathsf s}(X)=E_{21/22}^{(9)}(\mathsf s;X)\). Define the linking
coefficient
\begin{equation}
 \lambda(\mathsf s):=
 -\frac{E_{9,\mathsf s}(a(\mathsf s))
         \bigl(1-q\,a(\mathsf s)^3\bigr)}{a(\mathsf s)^{10}}.
 \label{eq:alpha-lambda-estado}
\end{equation}
All its arguments belong to the state before metrological recognition.
In particular, \(\lambda\) is not obtained by fitting a target decimal
string. Equation~\eqref{eq:alpha-lambda-estado} must be read with
its exact type: once the state's first representation has produced
\(a(\mathsf s)\), it fixes the unique extension within the geometric class
defined below. It is not a second independent selection of
\(a(\mathsf s)\), but a correlated analytic representation of the same
enriched state.

Indeed, for \(\mu\in F_{\mathsf s}\), set
\[
 E_{\mathsf s,\mu}(X)
 :=E_{9,\mathsf s}(X)+\mu\frac{X^{10}}{1-qX^3}.
\]
Since \(0<a(\mathsf s)<10^{-2}\) and \(qa(\mathsf s)^3<1\), we have
\begin{equation}
 E_{\mathsf s,\mu}(a(\mathsf s))=0
 \quad\Longleftrightarrow\quad
 \mu=\lambda(\mathsf s).
 \label{eq:alpha-unicidad-lambda-clase-geometrica}
\end{equation}
Thus the coefficient seed has no freedom once the state and the class
of \(q\)-geometric extensions are fixed. This uniqueness is relative to
the specified class; it does not assert that the order-nine jet alone selects
a tail from all possible analytic series.

The local action on blocks of three coefficients is
\begin{equation}
 \mathcal C_\alpha:F_{\mathsf s}^{3}\longrightarrow F_{\mathsf s}^{3},
 \qquad \mathcal C_\alpha(u,v,w)=q(u,v,w),
 \qquad c^{(0)}=(\lambda(\mathsf s),0,0).
 \label{eq:alpha-accion-coeficientes}
\end{equation}
Consequently, for every \(n\geq0\),
\begin{equation}
 b_{10+3n}=q^n\lambda(\mathsf s),\qquad
 b_{11+3n}=b_{12+3n}=0.
 \label{eq:alpha-recurrencia-coeficientes}
\end{equation}
This is the rule that computes every coefficient beyond degree nine.
No free parameter remains at each extension.

For \(M\geq0\), let
\begin{equation}
 E_{\mathsf s,M}(X):=E_{9,\mathsf s}(X)
 +\lambda(\mathsf s)\sum_{n=0}^{M}q^nX^{10+3n}.
 \label{eq:alpha-truncamientos-completos}
\end{equation}
The complete function is
\begin{equation}
 E_{\mathsf s,\infty}(X)
 :=E_{9,\mathsf s}(X)
 +\lambda(\mathsf s)\frac{X^{10}}{1-X^3/729}.
 \label{eq:alpha-funcion-completa-explicita}
\end{equation}

To formulate the projective system without implicit types, fix
\(d_M=12+3M\) and define the jet fibration
\begin{equation}
 \mathfrak J_M:=
 \coprod_{\mathsf s\in\mathcal X_\alpha^{\mathrm{enr}}}
 \{\mathsf s\}\times F_{\mathsf s}[X]/(X^{d_M+1}).
 \label{eq:alpha-espacio-jets-rev08}
\end{equation}
For \(M'\ge M\), restriction is
\begin{equation}
 \tau_{M',M}(\mathsf s,[P]_{d_{M'}})
 :=(\mathsf s,[P]_{d_M}).
 \label{eq:alpha-truncamiento-jets-rev08}
\end{equation}

\begin{proposition}[Naturality of the analytic recurrence]
\label{prop:alpha-naturalidad-recurrencia-explicita}
Let
\[
 \Gamma_M(\mathsf s)
 =\bigl(\mathsf s,[E_{\mathsf s,M}]_{d_M}\bigr).
\]
For \(M'\geq M\), truncation \(\tau_{M',M}\) satisfies
\begin{equation}
 \tau_{M',M}\circ\Gamma_{M'}=\Gamma_M,
 \qquad
 \Gamma_{E21}(\mathsf s)=\varprojlim_M\Gamma_M(\mathsf s).
 \label{eq:alpha-naturalidad-explicita}
\end{equation}
Projection onto the three new degrees is precisely
\(\mathcal C_\alpha^{M+1}(c^{(0)})\).
\end{proposition}

\begin{proof}
Equation~\eqref{eq:alpha-recurrencia-coeficientes} fixes the block of
degrees \(10+3n,11+3n,12+3n\). Passing from \(M\) to \(M+1\) adds only
the next block; truncating it recovers the previous polynomial exactly.
Compatibility for \(M'\geq M\) follows by composition. Thus the inverse
limit is not a section chosen from the bare jet, but the unique orbit
of \(c^{(0)}\) under \(\mathcal C_\alpha\).
\end{proof}

\subsubsection{Uniform convergence, a simple root, and rational intervals}
\label{subsec:alpha-cotas-completas-rev08}

The regional cylinders and periodic record provide bounds through
rational interval arithmetic. For each coordinate
\[
 c\in\{\pi_{\mathrm{HMT}},e_{\mathrm{HMT}},\varphi_{\mathrm{HMT}}\},
\]
let \(P_c\) be its certified nonadic prefix of one thousand digits. The datum
used is not the rational \(P_c\) treated as the complete value, but the
half-open cylinder and its outer closure
\[
 \mathcal C_c^{(1000)}=[P_c,P_c+\varepsilon),
 \qquad
 \overline{\mathcal C}_c^{(1000)}=[P_c,P_c+\varepsilon],
 \qquad \varepsilon=10^{-1000}.
\]
Directed arithmetic operates conservatively on these closures. If
\(P_\pi,P_e,P_\varphi\) denote the three prefixes, this gives
\begin{equation}
\begin{aligned}
 a_-&:=P_\pi+P_e-(P_\varphi+\varepsilon)-4-
       \kappa_{\mathrm{per}},\\
 a_+&:=(P_\pi+\varepsilon)+(P_e+\varepsilon)-P_\varphi-4-
       \kappa_{\mathrm{per}},\\
 A_{1000}&:=[a_-,a_+],
 \qquad a(\mathsf s)\in[a_-,a_+]=A_{1000},
 \qquad \operatorname{diam}A_{1000}=3\varepsilon.
\end{aligned}
\label{eq:alpha-propagacion-colas-rev08}
\end{equation}
The natural interval extension of
\eqref{eq:alpha-lambda-estado} then transports \(A_{1000}\), the cylinder
of \(\pi_{\mathrm{HMT}}\), and the directed rational bound for \(\delta\) to a
closed interval \(\Lambda_{1000}\ni\lambda(\mathsf s)\); none of the three
tails is frozen. In particular,
\begin{equation}
\begin{aligned}
 \frac{7297352569283800997285105472380662}{10^{36}}
 &<a(\mathsf s)<
 \frac{7297352569283800997285105472380663}{10^{36}},\\
 -8.711\mathbin{\times}10^{-6}
 &<\lambda(\mathsf s)<-8.709\mathbin{\times}10^{-6}.
\end{aligned}
 \label{eq:alpha-cotas-precoordenada-lambda}
\end{equation}
These inequalities use certified prefixes with their rational tail bound;
they receive no alpha table. The package verifier reproduces them
from the three nonadic outputs and the terminal chart.

Let
\[
 r:=\frac{10^{-6}}{729}<1.372\mathbin{\times}10^{-9}.
\]
For \(0\leq x\leq10^{-2}\), the tail beyond \(M\) satisfies
\begin{equation}
 \left|E_{\mathsf s,\infty}(x)-E_{\mathsf s,M}(x)\right|
 \leq
 8.711\mathbin{\times}10^{-26}\,\frac{r^{M+1}}{1-r}.
 \label{eq:alpha-cota-cola-uniforme}
\end{equation}
Moreover,
\begin{equation}
 \left|\frac{\mathrm d}{\mathrm dx}
  \left(\lambda\frac{x^{10}}{1-x^3/729}\right)\right|
 \leq 8.711\mathbin{\times}10^{-6}
 \left(\frac{10\,10^{-18}}{1-r}
 +\frac{3\,10^{-24}/729}{(1-r)^2}\right)
 <8.712\mathbin{\times}10^{-23}.
 \label{eq:alpha-cota-derivada-cola}
\end{equation}
The derivative tail beyond block \(M\) also admits the
depth-dependent estimate
\begin{equation}
 \sup_{0\le x\le10^{-2}}
 \left|E'_{\mathsf s,\infty}(x)-E'_{\mathsf s,M}(x)\right|
 \le 8.711\mathbin{\times}10^{-24}\,r^{M+1}
 \left(\frac{13+3M}{1-r}+\frac{3r}{(1-r)^2}\right)
 \xrightarrow[M\to\infty]{}0.
 \label{eq:alpha-cota-resto-derivadas-rev08}
\end{equation}

\begin{theorem}[Limit function and unique simple root of the correlated chart]
\label{thm:alpha-raiz-completa-explicita}
The sequence \(E_{\mathsf s,M}\) converges uniformly, together with its
derivatives, to \(E_{\mathsf s,\infty}\) on
\(\overline I_\alpha=[0,10^{-2}]\). We have
\begin{equation}
 E_{\mathsf s,\infty}(a(\mathsf s))=0,
 \qquad
 \inf_{x\in\overline I_\alpha}
 E'_{\mathsf s,\infty}(x)>6.239.
 \label{eq:alpha-raiz-y-derivada-completas}
\end{equation}
Consequently, \(a(\mathsf s)\) is the unique simple root of the complete
function in \(I_\alpha\).
\end{theorem}

\begin{proof}
The ratio of two successive terms in the tail is bounded by \(r<1\),
which gives \eqref{eq:alpha-cota-cola-uniforme}; termwise differentiation
and summation of the two geometric series give
\eqref{eq:alpha-cota-derivada-cola} and
\eqref{eq:alpha-cota-resto-derivadas-rev08}. The jet proof establishes
\(E'_{9,\mathsf s}>6.24\) on \(\overline I_\alpha\); hence the complete
perturbation leaves the derivative above \(6.239\).
Finally, substituting \eqref{eq:alpha-lambda-estado} into
\eqref{eq:alpha-funcion-completa-explicita} makes the two summands evaluated at
\(a(\mathsf s)\) cancel exactly. The function is strictly increasing;
its unique root is therefore simple and coincides with that pre-coordinate.
\end{proof}

To make identification effective at every depth, directed endpoint
evaluations first certify
\[
 E_{\mathsf s,\infty}(0)<0<
 E_{\mathsf s,\infty}(10^{-2}),
 \qquad
 E'_{\mathsf s,\infty}\geq\mu:=6239/1000>c:=6
 \quad\hbox{on }[0,10^{-2}].
\]
Continuity guarantees a root in the interval, and the second inequality
ensures uniqueness. Only after this existence, monotonicity, and the invariant
that the bracket contains the root have been established is the mean-value
theorem fallback applied. Start with \(D_0=[0,10^{-2}]\). If
\(D_j=[r_j,s_j]\), let \(m_j=(r_j+s_j)/2\), and let
\([F_j^-,F_j^+]\) be a directed rational enclosure of
\(E_{\mathsf s,\infty}(m_j)\), obtained by propagating the cylinders of
\(\pi_{\mathrm{HMT}},e_{\mathrm{HMT}},\varphi_{\mathrm{HMT}}\),
\(\delta\), and \(\lambda\), with
\begin{equation}
 0\leq F_j^+-F_j^-\leq \frac{c}{4}\,(s_j-r_j).
 \label{eq:alpha-anchura-evaluacion-rescate-rev08}
\end{equation}
When \(F_j^+<0\), retain the right half;
when \(F_j^->0\), retain the left half. If
\(F_j^-\leq0\leq F_j^+\), no attempt is made to decide whether the midpoint
is exactly the root. Let \(w_j=s_j-r_j\). The uniform bound
\(E'_{\mathsf s,\infty}\geq c=6\) and the mean value theorem imply that
every root contained in \(D_j\) belongs to
\begin{equation}
 D_{j+1}=D_j\cap
 \left[m_j-\frac{w_j}{4},
             m_j+\frac{w_j}{4}\right].
 \label{eq:alpha-rescate-derivada-rev08}
\end{equation}
Indeed, let \(\xi_j\) be the root, and let
\(y_j=E_{\mathsf s,\infty}(m_j)\in[F_j^-,F_j^+]\). Since the enclosure
contains zero and satisfies \eqref{eq:alpha-anchura-evaluacion-rescate-rev08},
we have \(|y_j|\leq cw_j/4\). The mean value theorem gives
\(|m_j-\xi_j|\leq |y_j|/c\leq w_j/4\), proving
\eqref{eq:alpha-rescate-derivada-rev08}. The argument includes the case
\(y_j=0\): the root then lies at the midpoint, but the algorithm need not
recognize that equality in order to retain it.

If the available enclosure does not satisfy
\eqref{eq:alpha-anchura-evaluacion-rescate-rev08}, the source cylinders are
first refined and the evaluation repeated; no sign is asserted and no
insufficient contraction is accepted. In each of the three branches, the new
interval retains the root and has diameter at most \(w_j/2\). For every
finite target depth, directed refinement makes it possible to satisfy the
width bound without an uncertified comparison of a real number with zero.
Monotonicity proves inductively that
\begin{equation}
 a(\mathsf s)\in D_{j+1}\subseteq D_j,\qquad
 \operatorname{diam}D_{j+1}\leq
 \tfrac12\operatorname{diam}D_j,
 \qquad \operatorname{diam}D_j\longrightarrow0,
 \label{eq:alpha-biseccion-racional}
\end{equation}
where the inequality is a contraction bound, not an imposed equality.
The executable certificate applies this rule at twenty-four levels in base
one thousand and verifies that the entire outer closure \(A_{1000}\), not
merely its lower endpoint, lies within each expressed cylinder. Its negative
tests reject an excessively wide enclosure, a reversed enclosure, a point
other than the midpoint, a domain without a sign change, and a degenerate
bracket. Coinductive extension allows one thousand to be replaced by any
required depth.

Specifically, for \(S_N=1000^N\), the verifier computes
\begin{equation}
 j_N^-:=\lfloor S_Na_-\rfloor,
 \qquad j_N^+:=\lfloor S_Na_+\rfloor
 \label{eq:alpha-indices-cilindro-dirigidos-rev08}
\end{equation}
and expresses the level only if \(j_N^-=j_N^+\). This equality, which includes
the upper endpoint of the outer closure, proves
\(A_{1000}\subset[j_N^-/S_N,(j_N^-+1)/S_N)\) and prevents an endpoint
lying exactly on the right boundary from being assigned to the cell on its left.

Let \(C_N^A(\mathsf s)\) be the integer cylinder in
\eqref{eq:alpha-cilindros-via-a-exactos}. Recursively choose
\(m(N)>m(N-1)\) so that
\(\operatorname{diam}D_{m(N)}<1000^{-N}/4\), and define
\begin{equation}
I_{B,N}:=D_{m(N)}\cap C_N^A(\mathsf s).
 \label{eq:alpha-intervalos-b-completos}
\end{equation}
Define
\(\ell_N:=\inf I_{B,N}\) and \(u_N:=\sup I_{B,N}\). Both are rational;
the interval is non-empty because both factors contain
\(a(\mathsf s)\). Since the cylinder is half-open, \(u_N\) need not
belong to \(I_{B,N}\). Evaluation at this endpoint is understood through
the continuous extension of \(E_{\mathsf s,\infty}\) to the outer closure. Thus,
\begin{equation}
 I_{B,N+1}\subseteq I_{B,N}\subseteq C_N^A(\mathsf s),
 \qquad \operatorname{diam}I_{B,N}\longrightarrow0,
 \qquad E_{\mathsf s,\infty}(\ell_N)\leq0\leq
 E_{\mathsf s,\infty}(u_N).
 \label{eq:alpha-inclusion-signos-completa}
\end{equation}
Bound~\eqref{eq:alpha-cota-cola-uniforme} allows a truncation to be chosen
that preserves the sign at every endpoint where the complete function is
non-zero. If an endpoint coincides with the root, the identity for the complete
function and the established membership are used; the same sign is not
attributed to its truncations. Indeed, the exact remainder is
\[
 E_{\mathsf s,\infty}(x)-E_{\mathsf s,M}(x)
 =\lambda(\mathsf s)x^{10}
   \frac{(qx^3)^{M+1}}{1-qx^3}.
\]
On the positive domain under consideration, where \(\lambda(\mathsf s)<0\),
evaluation at the root gives
\[
 E_{\mathsf s,M}(a(\mathsf s))
 =-\lambda(\mathsf s)a(\mathsf s)^{10}
   \frac{(qa(\mathsf s)^3)^{M+1}}{1-qa(\mathsf s)^3}>0
 \qquad(M<\infty).
\]
Accordingly, certification of the weak sign inequalities in
\eqref{eq:alpha-inclusion-signos-completa} concerns
\(E_{\mathsf s,\infty}\). Evaluation uses the complete function or the
truncation together with its signed remainder. This covers both the included
infimum and the potentially excluded supremum of the cylinder.

\Needspace{44\baselineskip}
\begin{proposition}[Compatibility squares and common cylinders]
\label{prop:alpha-cuadrados-compatibilidad}
Restriction of words and of analytic representations commutes:
\[
\begin{array}{ccc}
 \mathcal X_\alpha^{\mathrm{enr}}&
 \xrightarrow{\ \mathsf T_{\alpha,N'}\ }&\mathcal W_{N'}\\[1mm]
 \Big\Vert&&\Big\downarrow\rho_{N',N}\\[1mm]
 \mathcal X_\alpha^{\mathrm{enr}}&
 \xrightarrow{\ \mathsf T_{\alpha,N}\ }&\mathcal W_N
\end{array}
\qquad
\begin{array}{ccc}
 \mathcal X_\alpha^{\mathrm{enr}}&
 \xrightarrow{\ \Gamma_{M'}\ }&\mathfrak J_{M'}\\[1mm]
 \Big\Vert&&\Big\downarrow\tau_{M',M}\\[1mm]
 \mathcal X_\alpha^{\mathrm{enr}}&
 \xrightarrow{\ \Gamma_M\ }&\mathfrak J_M .
\end{array}
\]
The collection \(I_{B,N}\) provides the effective connection between the two
squares and explicitly satisfies
\(I_{B,N}\subseteq C_N^A\) for every \(N\).
\end{proposition}

\begin{proof}
The first square expresses compatibility of integer normalization; the second
is Proposition~\ref{prop:alpha-naturalidad-recurrencia-explicita}. The inclusion
follows from Definition~\eqref{eq:alpha-intervalos-b-completos},
and non-emptiness and the sign inequalities follow from
Theorem~\ref{thm:alpha-raiz-completa-explicita}. No step identifies the
root of a truncation with the root of an arbitrary tail.
\end{proof}

We now define, with the prefix reader made explicit,
\begin{equation}
 \alpha_B(\mathsf s):=
 \operatorname*{arg\,zero}_{x\in I_\alpha}E_{\mathsf s,\infty}(x),
 \qquad
 \operatorname{pref}_N(x):=
 \frac{\lfloor1000^Nx\rfloor}{1000^N},
 \qquad
 \alpha_{B,N}:=\operatorname{pref}_N(\alpha_B).
 \label{eq:alpha-definicion-b-y-prefijos-rev08}
\end{equation}

\begin{theorem}[Projective equality of the two correlated representations]
\label{prop:alpha-dos-limites}
For every \(N\), the integer route and the analytic route express the same prefix:
\begin{equation}
 \boxed{\alpha_{A,N}=\operatorname{pref}_N(\alpha_A)
 =\operatorname{pref}_N(\alpha_B)=\alpha_{B,N}.}
 \label{eq:alpha-dos-vias-cada-nivel}
\end{equation}
In the limit,
\begin{equation}
 \boxed{\alpha_A
 =\operatorname*{arg\,zero}_{x\in I_\alpha}
       E_{\mathsf s,\infty}(x)
 =\alpha_B=\alpha_{\mathrm{HMT}}.}
 \label{eq:alpha-dos-vias-limite-coinductivo}
\end{equation}
\end{theorem}

\begin{proof}
Integer normalization expresses exactly \(a(\mathsf s)\).
Theorem~\ref{thm:alpha-raiz-completa-explicita} proves that the analytic
chart represents the same point as its unique simple root. The half-open
cylinders \(C_N^A\) fix a canonical prefix, and
\(I_{B,N}\subseteq C_N^A\) contains \(\alpha_B=\alpha_A\). Their prefixes
therefore coincide for every \(N\); nesting and diameters tending to zero
give equality of the limits. Metrological recognition occurs
afterwards.
\end{proof}

The analytic function constitutes a correlated representation of the same
state, not a second derivation causally independent of the integer closure.
This precision does not weaken the proved equality: it identifies its
exact mechanism, prevents the jet from being assigned a tail it does not
determine, and avoids promoting a conventional value to a generator.
